%% UI Reference -- every control, metric, table column and formula in app.py.
%% Build:  cd paper && pdflatex ui_reference && pdflatex ui_reference
\documentclass[11pt]{article}

\usepackage[margin=1in]{geometry}
\usepackage{amsmath,amssymb}
\usepackage[hidelinks]{hyperref}
\usepackage{enumitem}
\usepackage{parskip}

\setlist[description]{style=unboxed, leftmargin=0pt, font=\normalfont\bfseries}

\title{Dynamic EMS Dispatch \& Station Placement Optimization\\[4pt]
\large UI Reference: Controls, Metrics and Formulas}
\author{}
\date{\today}

\begin{document}
\maketitle
\tableofcontents

\section*{How to read this document}
The application has a persistent sidebar and four tabs: \emph{Data}, \emph{Site
Scoring}, \emph{Station Placement} and \emph{Dispatch}. Each section below
matches one of those regions, in the order its controls and outputs appear on
screen. Every term is given a plain-language definition; where it is backed by
a formula, the formula follows immediately.

\section{Sidebar: Study Period and Response-Time Model}

\subsection{Study period}
\begin{description}
  \item[Quick range] A preset date window applied to the whole app: \emph{Last
    30 days}, \emph{Last 12 months} (default), \emph{Last 3 years},
    \emph{Entire dataset}, or \emph{Custom}. Every chart, table and model
    downstream is computed only on incidents inside the chosen window.
  \item[Date range (Custom)] Appears only when \emph{Quick range} is set to
    \emph{Custom}. A start and end calendar date bounding the incidents used.
  \item[Include trauma \& rescue calls] A toggle, off by default. When off,
    only \emph{core} EMS dispatches count as demand (calls whose type matches
    aid, medic, medical, EMS, AMR, nurseline, triaged, low acuity, overdose,
    cardiac or CPR). When on, \emph{trauma} calls are added as well: those
    typed by their cause rather than their medical nature (MVI/motor vehicle,
    rescue, scenes of violence, assault, casualty, drowning, extrication).
\end{description}

\subsection{Response-time model}
\begin{description}
  \item[Response-time limit (min)] The single number every optimisation model
    is steered by: the maximum time from alarm to on-scene arrival that counts
    as ``covered.'' Chosen from \{4, 5, 6, 7, 8, 9, 10, 12, 15\} minutes.
  \item[Implied reach] A read-only caption under the limit slider, converting
    the time limit into the straight-line distance it corresponds to under the
    current speed/circuity/turnout settings:
    \[
      r \;=\; \max(t_{\text{limit}} - \tau,\,0)\cdot\frac{v}{60\,\kappa}
    \]
    where $t_{\text{limit}}$ is the response-time limit (min), $\tau$ is
    turnout time (min), $v$ is effective speed (km/h) and $\kappa$ is the road
    circuity factor. $r$ is reported in km.
\end{description}

\subsubsection*{Model assumptions (expander)}
\begin{description}
  \item[Effective speed (km/h)] The assumed average road speed of a responding
    ambulance, $v$. Slider range 20--80, default \textbf{45}.
  \item[Road circuity factor] A multiplier $\kappa \geq 1$ that inflates
    straight-line (great-circle) distance into an estimate of actual
    street-network distance, since real roads are not straight lines. Slider
    range 1.0--1.8, default \textbf{1.30}.
  \item[Turnout time (min)] A fixed delay $\tau$ added to every response,
    representing the time between an alarm sounding and the crew actually
    leaving the station -- before any driving begins. Slider range 0--4,
    default \textbf{1.0}.
  \item[Travel-time formula] The equation these three constants feed, used by
    every model in the app to turn a straight-line distance into a response
    time:
    \[
      t_{ij} \;=\; \tau \;+\; \frac{60\,\kappa\, d_{ij}}{v}
    \]
    where $d_{ij}$ is the great-circle (Haversine) distance in km between
    point $i$ and point $j$, and $t_{ij}$ is the resulting travel time in
    minutes. The Haversine distance itself is
    \[
      a = \sin^2\!\Big(\tfrac{\Delta\phi}{2}\Big) + \cos\phi_1\cos\phi_2
          \sin^2\!\Big(\tfrac{\Delta\lambda}{2}\Big), \qquad
      d = 2R\arcsin\!\big(\sqrt{a}\big)
    \]
    with $R = 6371.0088$ km, $\phi$ latitude and $\lambda$ longitude in
    radians.
  \item[Charge within-zone travel] A toggle, \textbf{off} by default. Every
    demand zone is represented by a single point (its mean call location), so
    with the toggle off a station is charged only turnout time to reach the
    zone it sits in -- the standard covering-location simplification, and the
    one every reported result in the project's paper uses. Switched on, each
    zone's own mean call radius (how far its real incidents typically sit from
    that single point) is added to every arrival time there:
    \[
      t_{ij}' \;=\; t_{ij} \;+\; \frac{\kappa}{v}\cdot
        \text{spread\_km}_{j}\cdot 60
    \]
    This makes the covering model more conservative (more stations needed for
    the same limit) because it no longer treats intra-zone travel as free.
  \item[Zone materiality floor (\% of demand)] A threshold $f_{\min}$, default
    \textbf{0.1\%}, slider range 0--1\%. A demand zone is kept in the
    optimisation models only if its share of total period demand clears the
    floor:
    \[
      \text{demand\_share}_j \;=\; \frac{d_j}{\sum_k d_k} \;\geq\; f_{\min}
    \]
    where $d_j$ is zone $j$'s incident count. Zones below the floor remain
    visible in every table but are excluded from the covering and dispatch
    models, so a handful of stray mutual-aid calls in a far suburb cannot
    force a dedicated station of their own.
\end{description}

\section{Data Tab}

\subsection{Headline metrics}
\begin{description}
  \item[EMS incidents] Total incident count in the selected period.
  \item[Demand zones] Number of ZIP-code zones that clear the materiality
    floor and are therefore planned around.
  \item[Avg per day] Mean incidents per calendar day in the period:
    \[
      \text{per\_day} \;=\; \frac{\#\text{incidents}}{\#\text{distinct calendar days present}}
    \]
  \item[Busiest zone] The single ZIP zone with the highest incident count in
    the period.
  \item[Peak hour] The clock hour (00--23) with the most incidents summed
    across every day in the period (e.g.\ every 13:00--14:00 slot, not one
    specific date).
\end{description}

\subsection{Where demand concentrates}
\begin{description}
  \item[Busiest-zones chart and table] The top 15 ZIP zones ranked by incident
    count, each row showing its \textbf{Incidents} count and \textbf{Share}
    of all incidents in the period:
    \[
      \text{share\_pct}_j \;=\; 100\cdot\frac{\text{incidents}_j}{\sum_k \text{incidents}_k}
    \]
  \item[Cumulative-share caption] States the combined share of all incidents
    carried by just those top 15 zones -- the fact that demand is spatially
    concentrated, which is what justifies optimising station placement rather
    than spacing stations evenly.
\end{description}

\subsection{When load spikes}
\begin{description}
  \item[Calls by hour (chart)] Incident count for each of the 24 clock hours,
    summed over the whole period, alongside the average per day for that hour
    ($\text{avg\_per\_day}_h = \text{incidents}_h / \#\text{days}$).
  \item[Calls by weekday (chart)] Incident count for each of the 7 weekdays,
    summed over the whole period.
  \item[Hour $\times$ weekday heatmap (expander)] A 7-row $\times$ 24-column
    grid of incident counts, showing which specific hour/weekday combinations
    carry the true load peaks (a pattern the two one-dimensional charts above
    can hide, since a weekday peak and an hour peak need not coincide).
  \item[Call-type mix (expander)] The 15 most frequent call-type labels in the
    raw feed, each with its share of total incidents.
  \item[Daily trend (expander)] Incident count per calendar day across the
    period. The final day is flagged as \textbf{partial} and marked
    accordingly on the chart whenever the data stops before 23:00 local time
    that day, since the 911 feed is harvested live and the last day is usually
    cut off mid-way and not comparable to a full day.
\end{description}

\subsection{Demand weights fed into both optimisations}
\begin{description}
  \item[Weight per zone (table)] Every zone's raw incident count
    (\textbf{Incidents}) and its weight relative to the single busiest zone
    (\textbf{Weight}):
    \[
      \text{weight\_pct}_j \;=\; 100\cdot\frac{d_j}{\max_k d_k}
    \]
    This is the literal weight the covering model's demand-coverage
    constraints and objective use -- the quantity ``\% of demand covered''
    is always measured against, elsewhere in the app.
  \item[Demand across the city (map)] One dot per demand zone, placed at its
    \textbf{demand point} -- the mean coordinate of every real incident in
    that zone, \emph{not} the ZIP polygon's geometric centroid, so the
    candidate site sits where the calls actually happen:
    \[
      (\bar\phi_j,\bar\lambda_j) \;=\; \frac{1}{n_j}\sum_{k\in j}(\phi_k,\lambda_k)
    \]
    Shading encodes the zone's incident count.
\end{description}

\section{Site Scoring Tab (AHP)}
This tab runs Saaty's Analytic Hierarchy Process to rank candidate station
sites on three criteria, using pairwise judgements collected from each team
member.

\subsection{Pairwise judgements}
\begin{description}
  \item[Team members] A count, 1--8, of how many independent judgement sets
    to combine (default: the number of team members on record). Each member
    fills the comparison table independently; their matrices are aggregated
    into one set of weights before anything downstream uses them.
  \item[Saaty's 1--9 scale] The fixed intensity scale every pairwise judgement
    below is made on: 1 = equal importance, 3 = moderate, 5 = strong, 7 = very
    strong, 9 = extreme, with 2/4/6/8 as intermediate steps. A reciprocal value
    (e.g.\ 1/3, 1/5) expresses the same intensity in the opposite direction.
  \item[Pairwise comparison sliders] One slider per pair of the three criteria
    (\emph{Incident Density} vs \emph{Equity Gap}, \emph{Incident Density} vs
    \emph{Coverage Potential}, \emph{Equity Gap} vs \emph{Coverage
    Potential}), per member. The three criteria themselves:
    \begin{description}
      \item[Incident Density] Number of emergencies in the candidate zone.
      \item[Equity Gap] Distance (km) to the nearest station Seattle already
        operates; farther is read as higher unmet need.
      \item[Coverage Potential] How much total demand is reachable from that
        zone within the current response-time limit.
    \end{description}
    Each slider is read back underneath as a plain-English sentence (e.g.\
    ``Incident Density is moderately more important than Equity Gap''), so the
    number on the track is never the only account of what was chosen.
  \item[Reciprocal comparison matrix] What each member's slider values are
    assembled into, one $3\times3$ matrix $A^{(k)}$ per member $k$:
    \[
      A^{(k)}_{ij} = a^{(k)}_{ij}, \qquad A^{(k)}_{ji} = \frac{1}{a^{(k)}_{ij}},
      \qquad A^{(k)}_{ii} = 1
    \]
\end{description}

\subsection{Judgement consistency}
\begin{description}
  \item[Aggregated comparison matrix] The four (or $n$) members' matrices
    combined element-wise by geometric mean -- the only averaging rule that
    keeps the result reciprocal, so that if the members' views of $i$ vs $j$
    average to 3, their view of $j$ vs $i$ automatically averages to $1/3$:
    \[
      \bar A_{ij} \;=\; \left(\prod_{k=1}^{n} A^{(k)}_{ij}\right)^{1/n}
    \]
  \item[AHP weights (chart)] The priority vector $w$ actually used everywhere
    downstream: the principal right eigenvector of $\bar A$, normalised to
    sum to 1.
    \[
      \bar A\,w = \lambda_{\max}\,w, \qquad
      w \;\leftarrow\; \frac{|w|}{\sum_c |w_c|}
    \]
  \item[$\lambda_{\max}$] The principal eigenvalue of the aggregated matrix,
    computed as $\lambda_{\max} = \text{mean}_i\big((\bar A w)_i / w_i\big)$. A
    perfectly consistent $n\times n$ matrix has $\lambda_{\max} = n$ exactly.
  \item[CI (Consistency Index)] How far $\lambda_{\max}$ departs from the
    perfectly-consistent value $n$:
    \[
      \mathrm{CI} \;=\; \frac{\lambda_{\max} - n}{n - 1}
    \]
  \item[RI (Random Index)] Saaty's published average CI of a large sample of
    \emph{randomly generated} reciprocal matrices of the same order $n$; the
    yardstick CI is judged against. For $n=3$, $\mathrm{RI}_3 = 0.58$.
  \item[CR (Consistency Ratio)] The normalised measure actually used to accept
    or reject a set of judgements:
    \[
      \mathrm{CR} \;=\; \frac{\mathrm{CI}}{\mathrm{RI}_n}
    \]
    Judgements are considered logically consistent, and the derived weights
    trustworthy, when $\mathrm{CR} < 0.10$; the app shows a green confirmation
    message in that case and a red warning otherwise, naming the measured
    value against the 0.10 cutoff.
  \item[Per-member consistency (table)] The same CR computed on each member's
    own matrix in isolation, with a pass/fail \textbf{Passes} flag
    ($\mathrm{CR} < 0.10$). This answers a question the aggregate CR cannot: a
    geometric mean of several individually-inconsistent matrices can still
    come out looking consistent, so each member's own figure is shown
    separately.
\end{description}

\subsection{Candidate sites ranked by priority}
\begin{description}
  \item[Max-normalisation] Before combining, every criterion's raw values are
    scaled by their own column maximum, which is valid because all three
    criteria are \emph{benefit} measures on a ratio scale with a true zero:
    \[
      \hat x_{ic} \;=\; \frac{x_{ic}}{\max_i x_{ic}}
    \]
  \item[Priority score] The ranking chart and the \textbf{Priority} column in
    the table, a weighted sum of the three normalised criteria using the AHP
    weights $w$ derived above:
    \[
      s_i \;=\; \sum_{c} w_c\,\hat x_{ic}
    \]
  \item[Candidate ranking (table)] One row per candidate zone, sorted by
    priority score descending, showing the raw \textbf{Incidents}
    (Incident Density), \textbf{Nearest stn} distance in km (Equity Gap),
    \textbf{In reach} incident count (Coverage Potential, i.e.\ demand
    reachable from this zone within the response-time limit) and the final
    \textbf{Priority} score.
\end{description}
This priority score does not by itself decide how many stations open; see
the tie-break description under \S\ref{sec:placement}.

\section{Station Placement Tab}
\label{sec:placement}

\subsection{Set up the run}
\begin{description}
  \item[Candidate sites considered (top-N by priority)] A slider limiting the
    covering model's candidate pool to the top $N$ zones by AHP priority
    score, so a user can study the effect of a smaller candidate set. Defaults
    to using every material zone.
  \item[Objective] A two-way choice of which covering model to solve:
    \begin{description}
      \item[Mode A -- minimum stations for full coverage] (Set Covering,
        LSCP). Opens the fewest possible stations such that every demand zone
        is within the response-time limit.
      \item[Mode B -- best coverage for a fixed budget] (Maximal Covering,
        MCLP). Opens exactly a chosen number of stations, $P$, to cover as
        much demand as possible within the limit.
    \end{description}
  \item[Station budget $P$] Appears only under Mode B. The exact number of
    stations the model is allowed to open, 1 up to the candidate-pool size
    (capped at 30 in the slider).
  \item[Mode A formulation] The integer program actually solved, including a
    deliberately tiny AHP-derived tie-break term:
    \[
      \min \; \sum_i y_i \;+\; \varepsilon\sum_i \text{pref}_i\, y_i
      \quad\text{s.t.}\quad
      \sum_{i \in N(j)} y_i \ge 1 \;\; \forall j, \qquad y_i \in \{0,1\}
    \]
    where $N(j) = \{i : t_{ij} \le \text{limit}\}$ is the set of candidates
    that can reach zone $j$ in time, $\varepsilon = 0.5 / n_{\text{candidates}}$
    (small enough that it can never justify opening an extra station), and
    \[
      \text{pref}_i \;=\; \frac{\max_k s_k - s_i}{\max_k s_k - \min_k s_k}
    \]
    so that among several equally-small station sets the solver prefers the
    one favouring higher-priority sites from the AHP ranking.
  \item[Mode B formulation]
    \[
      \max \; \sum_j d_j\, z_j
      \quad\text{s.t.}\quad
      z_j \le \sum_{i \in N(j)} y_i \;\;\forall j, \qquad
      \sum_i y_i = P, \qquad y_i, z_j \in \{0,1\}
    \]
    where $d_j$ is zone $j$'s demand, so the objective maximises people
    reached rather than area or zone count.
\end{description}

\subsection{Result}
\begin{description}
  \item[Stations opened] The number of candidate sites the solver actually
    selected ($\sum_i y_i$).
  \item[Demand covered] Share of total period demand served by at least one
    opened station within the limit:
    \[
      \text{demand\_pct} \;=\; 100\cdot\frac{\sum_{j:\,\text{covered}} d_j}{\sum_j d_j}
    \]
  \item[Zones covered] Count of demand zones reached by at least one opened
    station within the limit, out of the total number of zones.
  \item[Solver status] The underlying ILP solver's own verdict, almost always
    \textbf{Optimal} on this problem size; other possible values are
    \textbf{Infeasible}, \textbf{Unbounded} or \textbf{Not Solved}.
  \item[Solve time] Wall-clock time the ILP solver spent on this run, shown in
    milliseconds below one second and in seconds above it.
  \item[Uncoverable-zones warning] Lists any zone that \emph{no} candidate
    site -- not just the opened ones -- can reach within the limit. Such zones
    are excluded from the covering constraints (otherwise the model would
    have no feasible solution at all) and are flagged so the gap is visible
    rather than silently absorbed.
  \item[Stations opened (table)] One row per opened station, with
    \textbf{Zone} (its ZIP) and \textbf{Serves}: the total demand of every
    zone whose nearest opened station, within the limit, is this one.
  \item[Per-zone outcome (table)] One row per demand zone: \textbf{Status}
    (\emph{Covered}/\emph{Not covered}), \textbf{From} (the opened station it
    is served by, if any) and \textbf{Response} (the resulting travel time in
    minutes to that nearest opened station).
  \item[Stations \& coverage (map)] Shaded circles show each opened station's
    response-time-limit reach; zone dots are green where covered and red
    where not.
\end{description}

\subsection{Sensitivity analysis}
\begin{description}
  \item[Run both sweeps (button)] Re-solves two families of covering problems
    over the current candidate pool: a \textbf{threshold sweep} (Mode A
    re-solved at every limit in $\{5, 8, 10\}$ minutes) and a \textbf{budget
    curve} (Mode B re-solved at every station budget from 1 up to a few past
    the current plan's station count, at the limit set in the sidebar).
  \item[Threshold-sweep table/chart] For each response-time limit: the
    minimum \textbf{Stations needed}, resulting \textbf{Demand covered} and
    \textbf{Zones covered}. Shows how tightening or relaxing the promise
    changes the station count a full-coverage plan requires.
  \item[Marginal coverage curve] Demand covered as a function of station
    budget $P$ under Mode B, plus the marginal gain of each additional
    station:
    \[
      \Delta_P \;=\; \text{demand\_pct}(P) - \text{demand\_pct}(P-1)
    \]
    Shows the diminishing return that justifies the budget-constrained (Mode
    B) framing: the first handful of stations buy most of the coverage.
\end{description}

\subsection{Robustness of the station count}
\begin{description}
  \item[Assumption grid] The Mode A problem re-solved at all 16 combinations
    of speed $v \in \{30, 35, 40, 45\}$ km/h and circuity
    $\kappa \in \{1.3, 1.4, 1.5, 1.6\}$, holding the candidate pool, AHP
    tie-break and limit fixed, so any spread in the answer is attributable to
    the travel-model assumption alone. Each grid cell also reports its
    implied straight-line reach, the same formula as the sidebar's
    \emph{Implied reach} caption:
    \[
      \text{reach\_km} \;=\; \max(\text{limit} - \tau,\,0)\cdot\frac{v}{60\,\kappa}
    \]
  \item[Stations needed (lo--hi)] The minimum and maximum station count found
    across all 16 grid cells at the current limit -- the honest range the
    headline number should be read within, rather than a single value.
  \item[Worst coverage] The lowest demand-covered percentage found in any
    single grid cell.
  \item[Still under Seattle's] How many fewer stations even the most
    pessimistic grid cell needs than the number of stations Seattle actually
    operates.
  \item[Verdict message] A green confirmation when every grid cell still
    needs fewer stations than Seattle's real network (meaning the placement
    advantage survives every tested assumption), or a red warning when the
    most pessimistic cell meets or exceeds Seattle's real count.
\end{description}

\subsection{Benchmark vs.\ real stations}
\begin{description}
  \item[Benchmark (button)] Scores both the optimised plan and the real
    Seattle Fire station network on identical travel assumptions, at every
    limit in $\{5, 8, 10\}$ minutes, using each network's own fixed set of
    stations (\texttt{evaluate\_fixed\_stations} for Seattle's, the Mode A
    solve for the optimised plan).
  \item[Comparison chart/table] \textbf{Stations}, \textbf{Demand covered} and
    \textbf{Zones covered} for \emph{Optimised placement} side by side with
    \emph{Seattle's real stations}, at each limit.
  \item[Headline ratio] At the sidebar's current limit, the real network's
    station count divided by the optimised plan's station count -- how many
    times fewer sites the optimised plan needs for the same coverage. The
    accompanying caption notes two caveats: Seattle's stations are
    multi-mission fire stations (also carrying suppression, hazmat and
    technical rescue), so the comparison is not perfectly like-for-like; and
    the optimised count is itself a function of the current travel
    assumptions (see \emph{Robustness}, above).
\end{description}

\section{Dispatch Tab}
This tab requires a solved covering result from \emph{Station Placement};
its opened stations become the ambulance bases dispatched from here.

\subsection{Choose the scenario}
\begin{description}
  \item[Peak hour of simultaneous emergencies] A dropdown of the 25 busiest
    individual clock hours in the selected period (by real incident count),
    used as the demand scenario for the dispatch model.
  \item[Busiest hours (table)] Those same candidate hours, with their
    timestamp, weekday and incident count.
  \item[Ready units per station] A slider, 1--12, for how many ambulances are
    presumed ready at each opened station. The resulting fleet size and the
    number of units that \emph{would} be needed to field every call that hour
    are both shown:
    \[
      \text{fleet} \;=\; \text{units}\times n_{\text{stations}}, \qquad
      \text{need} \;=\; \left\lceil \frac{n_{\text{calls}}}{n_{\text{stations}}} \right\rceil
    \]
    Defaulting the slider to \texttt{need} (capped at 12) avoids starving the
    model with far fewer ambulances than the hour actually requires.
\end{description}

\subsection{Scenario metrics}
\begin{description}
  \item[Incidents this hour] Real call count inside the chosen clock hour.
  \item[Ambulances available] Fleet size, $\text{units}\times n_{\text{stations}}$.
  \item[Capacity gap] Calls with no ambulance left to send \emph{before} the
    model is even solved: $\max(n_{\text{calls}} - n_{\text{ambulances}}, 0)$.
  \item[Response-time limit] The limit currently set in the sidebar, used
    here only to judge the dispatch results against -- deliberately
    independent of the (possibly different) limit the stations were
    \emph{placed} under in the Station Placement tab, so ``place at 8
    minutes, judge against 6'' is a question the app can actually answer.
\end{description}

\subsection{Dispatch result}
\begin{description}
  \item[Assignment LP] The model solved on pressing \emph{Solve the
    assignment model}: minimise total travel time across all ambulance/
    incident pairings, one ambulance per incident:
    \[
      \min \sum_{i,j} c_{ij}\,x_{ij}
      \quad\text{s.t.}\quad
      \sum_i x_{ij} = 1 \;\;\forall j, \qquad
      \sum_j x_{ij} \le 1 \;\;\forall i, \qquad
      x_{ij} \ge 0
    \]
    where $c_{ij}$ is the travel time from ambulance base $i$ to incident
    $j$ (the same travel-time formula as the sidebar). The constraint matrix
    is totally unimodular, so this LP relaxation is already guaranteed an
    integral optimum -- no binary variables are needed. Unequal ambulance and
    incident counts are first balanced by padding the cost matrix with
    zero-cost dummy rows (fewer ambulances than calls: dummy ambulances
    absorb the calls that cannot be served) or dummy columns (more ambulances
    than calls: dummy incidents absorb idle units).
  \item[Total travel time] Sum of travel time over every call this plan
    actually served.
  \item[Mean response] Average travel time over served calls.
  \item[Worst response] Maximum travel time over served calls.
  \item[Served / unserved] Count of calls assigned a real ambulance versus
    calls left with none (shown in the dispatch plan table as
    \emph{Unserved}).
  \item[Idle units] Ambulances assigned no call this hour:
    $\max(n_{\text{ambulances}} - n_{\text{served}},\,0)$.
  \item[Dispatch plan (table)] One row per incident: the responding
    \textbf{Unit} (or \emph{Unserved} if none was available), the base's
    \textbf{From} zone, and the resulting \textbf{Response} time (an em dash
    where there is none).
\end{description}

\subsection{Benchmark against greedy dispatch}
\begin{description}
  \item[Greedy nearest-free-unit baseline] What dispatch looks like without
    optimisation: calls are handled in arrival order, each sent the nearest
    still-free unit.
    \[
      i^*_j \;=\; \arg\min_{i \in \text{free}} c_{ij}, \qquad
      \text{free} \;\leftarrow\; \text{free}\setminus\{i^*_j\}
    \]
    This is the benchmark the optimal assignment is meant to beat: a locally
    nearest choice can strand a later call with a very long drive.
  \item[Like-for-like restriction] Both the optimal LP total and the greedy
    total are computed over the \emph{same} $k = \min(n_{\text{ambulances}},
    n_{\text{incidents}})$ calls, in arrival order -- the workload a
    dispatcher actually has to clear. Without this restriction, an
    unconstrained LP would quietly serve only the cheapest calls when the
    fleet is too small, overstating its advantage; greedy has no such choice,
    since it must take calls as they arrive.
  \item[Saved vs.\ nearest-unit-first] The routing advantage of the optimal
    matching over the greedy baseline, on that same restricted workload:
    \[
      \text{improvement\_pct} \;=\; 100\cdot
      \frac{\text{greedy\_total} - \text{optimal\_total}}{\text{greedy\_total}}
    \]
  \item[Response distribution (chart)] Histogram of response times across
    served calls this hour, against the current response-time limit.
  \item[Dispatch comparison (chart)] Total travel time, optimal LP vs.\
    greedy, over the same $k$ calls.
\end{description}

\subsection{Supporting detail}
\begin{description}
  \item[Cost matrix (expander)] The real ambulance-by-incident travel-time
    matrix in minutes, and, whenever the counts were unequal, the
    \textbf{Balanced matrix} padded with zero-cost dummy ambulances or dummy
    incidents to make the assignment problem square.
  \item[Dispatch plan (map)] A line from each ambulance's base to the
    incident it was sent to; red dots mark calls no unit could reach.
\end{description}

\end{document}
